% !PRINTED 574
\weberpage{574}
Lassen wir in $\Phi(t)$ und $\Phi_0(t)$ Glieder weg, deren Coëfficienten durch $p$ theilbar sind, so ist der Grad von $\Phi(t)$ in Bezug auf $t$ um $f$ grösser als der Grad von $\Phi_0(t)$.

$P(t)$ ist eine ganze Function in $R$ der Variablen $t, t_1, t_2, \ldots, t_n$, die durch die Substitution $t=\tau$ in ein durch $\mathfrak p$ theilbares Functional übergeht, die natürlich nicht mehr in $R$ enthalten ist.

Die Bedeutung dieser Form $P(t)$ tritt nun noch deutlicher hervor, wenn wir den Satz beweisen:

4. \gesperrt{Der Primfactor $\mathfrak p$ ist der grösste gemeinschaftliche Theiler von $p$ und $P(\tau)$.}

Dazu haben wir nur nachzuweisen, dass, wenn $p$ durch $\mathfrak p \mathfrak p_1$ theilbar ist, wo $\mathfrak p, \mathfrak p_1$ zwei gleiche oder verschiedene Primfactoren sind, $P(\tau)$ zwar durch $\mathfrak p$, nicht aber durch $\mathfrak p \mathfrak p_1$ theilbar ist.

Nach \S.\,143, 3. existirt immer eine \gesperrt{Zahl} $\xi$ in $\mathfrak{o}$, die zwar durch $\mathfrak p$, aber nicht durch $\mathfrak p \mathfrak p_1$ theilbar ist. Diese Zahl wird die Form haben
\weq{16}{\xi = a_1 \omega_1 + a_2 \omega_2 + \cdots + a_n \omega_n,}
worin die $a_1, a_2, \ldots, a_n$ ganze rationale Zahlen sind, und geht also aus der Form $\tau$ hervor durch die Substitution
\weq{17}{(t_1, t_2, \ldots, t_n) = (a_1, a_2, \ldots, a_n).}
Wenn wir dieselbe Substitution in den in (9) definirten Formen $\tau_r$ machen, so geht $\tau_r$ in eine Zahl $\xi_r$ über, die nach dem \gesperrt{Fermat}'schen Satze der Congruenz
\[ \xi^{p^r} \equiv \xi_r \pmod{\mathfrak p} \]
genügt und also sicher auch durch $\mathfrak p$ theilbar ist. Wenn wir daher in der Form
\[ \Pi(t) = (t-\tau)(t-\tau_1) \cdots (t - \tau_{f-1}) \]
$\tau_r + \xi_r$ an Stelle von $\tau_r$ setzen, so bleibt diese Form mit sich selbst nach dem Modul $\mathfrak p$ congruent. Diese Substitution kommt aber darauf hinaus, dass wir $t_1+a_1, t_2+a_2, \ldots, t_n+a_n$ an Stelle von $t_1, t_2, \ldots, t_n$ substituiren, und wir erhalten demnach
\[ \Pi(t, t_1+a_1, \ldots, t_n+a_n) \equiv \Pi(t, t_1, \ldots, t_n) \pmod{\mathfrak p}, \]
% [sic: Weber prints ``nnd'' for ``und''.]
nnd wenn wir $\Pi$ durch die congruente Form $P$ ersetzen
\[ P(t, t_1+a_1, \ldots, t_n+a_n) \equiv P(t, t_1, \ldots, t_n) \pmod{\mathfrak p}. \]
